% Part: second-order-logic
% Chapter: sol-and-set-theory
% Section: introduction

\documentclass[../../../include/open-logic-section]{subfiles}

\begin{document}

\olfileid{sol}{set}{int}

\olsection{પ્રસ્તાવના}

દ્વિતીય-ક્રમ તર્કશાસ્ત્રમાં વ્યાપકક્ષેત્રના ઉપગણો અને વિધેયો પર
પરિમાણન કરી શકાતું હોવાથી, તેમાં ઓછામાં ઓછું થોડું ગણસિદ્ધાંતનું
કામ કરી શકાય એવી અપેક્ષા સ્વાભાવિક છે. ``કામ કરી શકાય'' તેનો અહીં
અર્થ એ છે કે ગણસિદ્ધાંતના ગુણધર્મો અને વિધાનો દ્વિતીય-ક્રમ
તર્કશાસ્ત્રમાં વ્યક્ત કરી શકાય છે; તે પણ ગણો માટેની કોઈ ખાસ
અતાર્કિક શબ્દાવલી વિના (દાખલા તરીકે, ગણસિદ્ધાંતના સભ્યતા
!!{predicate} વિના). દાખલા તરીકે, આપણે ગણોનો સંયોગ અને છેદ તથા
ઉપગણ-સંબંધ વ્યાખ્યાયિત કરી શકીએ છીએ; ગણોનાં કદની સરખામણી પણ કરી
શકીએ છીએ અને કૅન્ટૉરના પ્રમેય જેવાં પરિણામો જણાવી શકીએ છીએ.

\end{document}
